\chapter{Empirical Application: Tablet Spectroscopy}
\label{chap:main-empirical}

The simulations distinguish slope recovery, prediction, and numerical
stability because the generating model is known.  This application examines
prediction and stability in pharmaceutical-tablet data.  The true slope is
unknown, so slope ISE cannot be assessed. Formal inference would additionally
require feasible calibration, whose performance was not evaluated in
Chapter~\ref{chap:main-inference}.  The detailed analysis is
preserved in Appendix~\ref{chap:spectroscopy-application}, with the complete
protocol and diagnostics in Appendix~\ref{app:spectroscopy-supplement}.

\section{Data and evaluation}
\label{main:sec:empirical-design}

The 2002 pharmaceutical-tablet NIR Shoot-out contains near-infrared spectra
from two instruments \parencite{hopkins_2003,eigenvector_tablets}.  Its 155
calibration, 40 validation, and 460 test tablets total 655, although the
catalogue states 654.  All complete records are retained.  The instruments
measured the same tablets, so their observations are paired.  The primary
analysis uses instrument~1 and all 650 wavelengths from 600 to 1898~nm in
2-nm steps.  Assay, the third response column, is the sole response; its
physical unit is unspecified, and errors use the supplied assay units.

Figure~\ref{main:fig:empirical-overview} shows the functional predictors and
the response distributions.  The spectra share broad shapes across the
development and benchmark samples, with a less regular region at the
long-wavelength end.  The comparison therefore asks how effectively each
procedure extracts predictive information from strongly related spectral
coordinates.  The visible tail motivates the prespecified trimming
sensitivity; its appearance alone does not establish that it is noise.

\begin{figure}[htbp]
    \centering
    \includegraphics[width=0.82\textwidth]{tablet_spectra_assay_overview.pdf}
    \caption[Tablet spectra and assay distributions]{Instrument-1 spectra and
    assay distributions.  The spectral panel shows the median and
    10th--90th percentile band for development and benchmark samples.  The
    dashed line marks 1798~nm, the boundary used only in the trimming
    sensitivity.  The response panel retains the three supplied splits.}
    \label{main:fig:empirical-overview}
\end{figure}

The validation split has a particularly narrow assay range, 188.2--199.5;
267 of the 460 benchmark responses lie outside it.  Its response standard
deviation is 2.393, compared with 21.980 in calibration and 15.763 in the
benchmark.  Consequently, validation-set $R^2$ has a much smaller
denominator and is not directly comparable with benchmark $R^2$.

Calibration and validation form a development sample of 195 tablets; the
460 test tablets form the designated benchmark.  The final evaluation rules
were fixed before fitting the four benchmark models.  However, benchmark
responses and preliminary ordinary PCR/PLS results had been inspected during
dataset selection.  This is therefore a designated benchmark, not a
prospectively blind validation of future performance.

Every training fold supplies its own wavelength and response means, which
are then applied unchanged to its validation observations.  The means are
recalculated when the training sample changes, including for the final fit
to all development tablets.  Benchmark observations never contribute to
these fitted means.  There is no additional
preprocessing in the primary analysis.  The equally weighted inner product
from Chapter~\ref{chap:main-methods} gives predictions
$\widehat Y(x)=\bar Y_{\mathrm{train}}+
(x-\bar X_{\mathrm{train}})^\top\widehat\beta/650$.
The centred development covariance has rank at most 194, leaving at least
456 unidentified coordinate directions.

Each method retains its own tuning rule.  Raw and Arnoldi FPLS independently
minimise prediction error over at most 70 requested dimensions using the
same five inner folds.  Raw uses unscaled covariance powers and
response-space normal equations; Arnoldi uses twice reorthogonalised
Gram--Schmidt and response-space QR.  FPCR selects its cut-off by
response-space GCV.  CG--FPLS uses the estimation study's conjugate-residual
recurrence and discrepancy rule, with $\tau=1.01$, $\delta=0.10$, and its
internal moment-GCV variance pilot.  These are comparisons of complete
procedures; component counts are not common effective degrees of freedom.

Development evaluation uses five repeats of five-fold outer
cross-validation, stratified by source split and balanced over assay.
Every outer fold contains 31 calibration and eight validation tablets;
tuning occurs within the remaining observations.  Each repeat provides one
out-of-fold summary.  The 25 overlapping fits are not independent
experiments.  Each final model is then fitted to all 195 development
tablets and predicts the 460 benchmark tablets once.

RMSE is the primary measure.  Uncertainty uses 10,000 paired iid bootstrap
resamples, with identical tablet indices across methods and fixed fitted
models.  Circular moving-block resamples of lengths 5, 10, and 20 assess
sensitivity to row ordering.  These intervals are conditional on the fitted
models and supplied benchmark: they omit development-sample, tuning, and
preprocessing uncertainty.

\section{Prediction and numerical stability}
\label{main:sec:empirical-results}

All four procedures give finite predictions and explain substantial assay
variation (Table~\ref{main:tab:empirical-results}).  Raw and Arnoldi have the
smallest benchmark RMSEs, followed by FPCR and CG--FPLS.

\begin{table}[htbp]
    \centering
    \small
    \caption[Primary tablet benchmark results]{Instrument-1 benchmark results
    ($n=460$).  Brackets give individual 95\% percentile intervals from
    10,000 paired iid bootstrap resamples.  Bias is prediction minus
    observation; $m$ is the selected dimension or stopping index.}
    \label{main:tab:empirical-results}
    \begin{tabular}{@{}lrrrrr@{}}
        \toprule
        Method & $m$ & RMSE [95\% interval] & MAE & Bias & $R^2$ \\
        \midrule
        CG--FPLS     & 3  & 5.293 [4.367, 6.477] & 3.673 &  0.177 & 0.887 \\
        Raw FPLS     & 55 & 4.607 [3.645, 5.785] & 3.036 & -0.303 & 0.914 \\
        Arnoldi FPLS & 5  & 4.643 [3.693, 5.805] & 3.086 & -0.301 & 0.913 \\
        FPCR         & 53 & 5.185 [4.299, 6.255] & 3.547 & -0.628 & 0.892 \\
        \bottomrule
    \end{tabular}
\end{table}

The methods share several difficult tablets
(Figure~\ref{main:fig:empirical-predictions}).  Raw and Arnoldi predictions
have correlation 0.999934 and root mean squared difference 0.174 assay units.
Their paired RMSE difference, Raw minus Arnoldi, is $-0.036$, with iid
interval $[-0.057,-0.017]$ and block-length-10 interval
$[-0.062,-0.014]$.  The difference is only 0.78\% of Arnoldi's RMSE;
its practical importance cannot be judged without an assay scale and
application-specific tolerance.  The narrow paired interval reflects close
agreement between the predictions and does not establish general superiority
of the raw representation.  CG-minus-Arnoldi and FPCR-minus-Arnoldi RMSEs
are 0.650 and 0.542, respectively; their iid and block-length-10 intervals
remain above zero.

\begin{figure}[htbp]
    \centering
    \includegraphics[width=0.82\textwidth]{tablet_benchmark_observed_vs_predicted.pdf}
    \caption[Observed and predicted tablet assay]{Observed and predicted assay
    for the 460 benchmark tablets.  The diagonal denotes perfect prediction.
    Several large errors are shared across methods; all tablets remain in
    the comparison.}
    \label{main:fig:empirical-predictions}
\end{figure}

The leading pair differs much more in numerical stability.  Arnoldi selects
four to six components across the 25 development fits, with median five and
no numerical flags.  Raw selects four to 70, with median 44, reaches the
cap twice, and is flagged in 24 of 25 fits.  Across the five repeats,
RMSE has mean and standard deviation $5.082\pm0.517$ for Raw versus
$4.917\pm0.131$ for Arnoldi.  One poor raw repeat contributes substantially,
and contiguous folds reverse their order, so these summaries do not
establish an independent performance ranking.

The final raw fit selects nominal dimension 55 and is flagged.  Its
normalised power-basis and response-design condition numbers reach the
reporting ceiling $4.50\times10^{15}$, not an exact measured value.  The
$1/\sqrt{\epsilon}$ threshold is crossed by dimension six.  Arnoldi selects
five components, with basis condition number one, design condition number
11.95, and orthogonality defect $7.4\times10^{-16}$.  Nearly identical
predictions therefore coexist with sharply different coordinate stability.

CG--FPLS stops at three components in every development and final fit.
Its larger error accompanies strong early regularisation, but this
comparison does not isolate the effect of extra CG iterations or show that
CG regularisation is generally inferior.  FPCR selects seven to 57
components in development and 53 in the final fit.

\section{Sensitivity and interpretation}
\label{main:sec:empirical-scope}

Prespecified sensitivity analyses change the detailed ranking.  Trimming
instrument~1 at 1798~nm brings all four RMSEs into the range 4.63--4.77,
with Arnoldi lowest.  Savitzky--Golay smoothing favours Raw at 4.66 but does
not improve every method.  Instrument~2 increases all four RMSEs; because
it measures the same tablets, this is an instrument sensitivity analysis,
not external validation or calibration transfer.

These findings neither identify a chemically causal slope nor establish
that visible spectral spikes caused particular errors.  Without batch,
acquisition-order, or laboratory metadata, residual patterns cannot be
assigned a physical explanation.  Within one response and one dataset, the
application supports the distinction established in
Chapter~\ref{chap:main-simulation}: useful prediction can coexist with
unreliable raw coordinates, and a numerically stable representation need
not change the predictions substantially.
